\documentclass[11pt]{article}
\usepackage[a4paper,margin=25mm]{geometry}
\usepackage{amsmath,amssymb,mathtools}
\usepackage{fontspec,unicode-math}
\setmainfont{Latin Modern Roman}
\setsansfont{Latin Modern Sans}
\setmonofont{Latin Modern Mono}
\setmathfont{Latin Modern Math}
\usepackage{xurl,hyperref}
\hypersetup{colorlinks=true,urlcolor=blue,linkcolor=blue}
\newcommand{\sref}[2]{\hyperlink{source:#1:#2}{[S#2]}}
\newcommand{\cataloguescope}[1]{\paragraph{Mathematical question.}#1}
\begin{document}
\section{Vandiver's conjecture}\label{statement}

\subsection{Definitions and mathematical statement}

\paragraph{Shared notation.}$\mathbb N=\{1,2,\ldots\}$, and $\mathbb Z,\mathbb Q,\mathbb R,\mathbb C$ denote the integers, rationals, reals and complex numbers. Any other conventions are specified locally.


For each prime $p$, set $\zeta_p=\exp(2\pi i/p)$ and $K_p^+=\mathbb Q(\zeta_p+\zeta_p^{-1})=\mathbb Q(\zeta_p)\cap\mathbb R$. This is the maximal real subfield of the $p$th cyclotomic field. Its ideal class group is the group of nonzero fractional ideals of its ring of integers, modulo principal fractional ideals; its order $h_p^+$ is the class number.
\paragraph{Statement recorded in the supplied source.}
Is $p\nmid h_p^+$ for every prime $p$? Equivalently, is the $p$-primary subgroup of the ideal class group of $K_p^+$ trivial for every prime $p$? \sref{statement}{2}

\subsection{Short English statement}
Does a prime never divide the class number of the maximal real subfield of its cyclotomic field?
\subsection{Sources}
\begin{description}
\item[\hypertarget{source:statement:1}{S1}] ULAM.AI and the UnsolvedMath Contributors, UnsolvedMath: A Curated Collection of Open Mathematics Problems, catalogue entry. CC BY 4.0. \url{https://huggingface.co/datasets/ulamai/UnsolvedMath}.
\item[\hypertarget{source:statement:2}{S2}] Christian Wuthrich, \emph{Open problems in number theory}, lecture slides (December 2011), slide headed ``Vandiver's conjecture, PDF page 31. \url{https://www.maths.nottingham.ac.uk/plp/pmzcw/download/open.pdf}.
\end{description}
\label{end:statement}
\end{document}
